%% 第一章 绪论

\chapter{绪论}
\chaptermark{绪论}

\section{研究背景与意义}

载人航天工程代表着国家的综合科技实力，而航天员的生命安全始终是整个系统设计的核心目标。从发射准备、点火升空到进入轨道前的全过程中，航天器随时可能遭遇复杂环境变化或突发故障风险。逃逸系统作为载人飞船的关键安全设施，在出现火箭爆炸、推力异常或姿态失控等险情时，能快速将返回舱与故障火箭分离，引导航天员进入可控着陆轨迹。这套瞬时脱险机制，为飞行任务增加了重要的安全冗余。

新一轮深空探测让载人登月再次成为国际航天领域的热点。与近地轨道任务相比，登月任务的轨道更复杂，能量需求更大，动力学环境也更为严苛，这对逃逸系统的可靠性提出了更高标准。我国新一代梦舟载人飞船要同时服务于近地空间站运营和载人探月两类任务——这两类任务的发射流程、工作时序和飞行弹道差异很大，逃逸系统需要具备较强的兼容性。特别是执行登月任务时，长征十号火箭的体量和潜在爆炸当量都大幅增加，而且濒海发射场的气象条件比内陆复杂得多。这些都要求逃逸系统具备更强的加速逃逸能力、更宽的适应边界，以及更精准的落点控制~\cite{田林2023}。

载人逃逸任务剖面设计是整个逃逸任务的第一步。从航天员进入飞船，到飞船进入轨道或安全着陆，每个时间段和飞行状态都需要考虑用什么样的逃逸方案，才能把航天员从危险环境带到安全区。整个逃逸过程的时间轴和弹道方案，就构成了载人逃逸任务剖面。从点火前待发、穿越大气层、接近轨道到入轨后的整个上升段，各阶段的飞行环境、气动载荷、结构状态和控制能力都不同，相应地，逃逸方式、分离策略和再入路径也要随之调整。设计逃逸方案时，必须保证无论故障在哪个时刻发生，飞船都能进入可控的逃逸弹道，最终实现可搜救着陆。

开展逃逸任务仿真研究，可以为载人飞船的逃逸动力学建模提供系统方法，丰富载人航天应急逃逸领域的理论积累；能推动六自由度动力学模型在复杂任务剖面设计中的应用；也能为后续深空任务的应急策略研究提供技术支持，进而提升飞船系统的整体可靠性。

本课题选取美国猎户座飞船作为参考对象，对发射逃逸阶段的相关问题进行初步探索。猎户座飞船与我国正在研制的梦舟新一代载人飞船属于同类航天器，在逃逸系统的设计思路上有一定共通之处。本次研究主要围绕发射逃逸阶段的六自由度动力学建模、蒙特卡洛仿真和风险评估这几个方向展开，希望能通过这项工作加深对载人飞船逃逸系统的理解，为今后相关领域的学习打下基础。

\section{国内外研究现状}

航天器的逃逸救援任务，按照飞行高度可以分为两个阶段：大气层内和大气层外高空段。在大气层内，逃逸塔是主要的救生装备——凭借高推重比，逃逸塔能在几秒内把返回舱拉离险境。当飞船接近入轨高度时，逃逸塔完成使命被抛弃，后续的应急分离与撤离主要依靠飞船自身的轨控发动机实现~\cite{张智2004}。

\subsection{国际研究现状}

美国在逃逸系统领域积累了丰富的工程经验~\cite{王鹏程2023}。阿波罗计划期间（1961--1972），美国完成了从发射台到返回舱的全流程逃逸设计，核心是发射逃逸系统和登月舱救生艇方案~\cite{McCarthy1968}。在阿波罗13号那次著名的危机中，登月舱被临时用作救生艇，帮助航天员安全返回地球。阿波罗飞船的逃逸剖面分为六大段：待发段、低空中止段、高空中止段、二级故障中止段和地月转移段。如果发射前5分钟待发阶段出现危险，逃逸系统点火将指令舱拉到1200米高空后抛塔开伞；要是30公里以下低空最大动压区发生故障，双发动机会协同控制指令舱转向，把航天员过载限制在12G以内；地月转移及登月段，服务舱失效时登月舱可以重构生命保障系统，为航天员提供支持。

美国猎户座飞船的发射逃逸系统沿用了阿波罗的基本方案，但做了不少改进。NASA先后开展了Pad Abort-1（PA-1，2010年）和Ascent Abort-2（AA-2，2019年）两次飞行试验，分别验证了零高度和最大动压条件下的逃逸能力~\cite{McCauley2015,Litton2008,Proud2009,徐丹丹2018}。

SpaceX的载人龙飞船另辟蹊径，采用集成在飞船侧壁的SuperDraco液体发动机实现推式逃逸，直接取消了传统的逃逸塔。龙飞船的逃逸剖面分为五个阶段：待发段遇到危险时，发动机可以在1.2秒内点火，把飞船推离危险区；低空最大动压区故障时，航天员过载控制在4.5G以内。这个方案消除了抛塔带来的额外死重和故障节点，同时发动机入轨后还能兼作轨控动力使用。

苏联/俄罗斯的联盟号飞船采用三舱构型，返回舱放在中段，逃逸塔装在整流罩顶部。联盟号的逃逸剖面分成7个阶段，覆盖从待发段到着陆段的全任务周期。2018年联盟MS-10任务发射失败时，飞船在约50公里高空启动了高空逃逸发动机，航天员经历6.7G过载后最终安全返回。这次发射虽然失败了，却验证了联盟号逃逸阶段时序划分和逃逸手段的合理性。

\subsection{国内研究现状}

我国载人航天科研人员在借鉴国外技术成果的基础上，独立突破了多项关键技术。神舟飞船的应急救生系统是按飞行阶段划分的，主要分为待发段、发射段大气层内、发射段大气层外、轨道运行段、返回段和着陆段这几个阶段~\cite{李颐黎2022}。1998年，我国完成了首次零高度逃逸飞行试验，验证了神舟飞船的应急救生能力。另外，为了模拟长征二号F火箭在最大动压附近的故障情况，我国还在1996年完成了最大动压滑轨试验，用火箭撬做了3次栅格翼释放展开试验，测试气动力对逃逸飞行器结构的影响。

我国的梦舟飞船要同时服务于载人登月和空间站任务，在逃逸任务剖面设计上和以往有很大不同。它把载人逃逸任务划分为两个阶段：大气层内用逃逸塔方案，大气层外由飞船自身承担逃逸任务，实现整船逃逸~\cite{田林2023}。2025年6月17日，梦舟飞船成功完成了零高度逃逸飞行试验。这次试验全面验证了逃逸救生分系统及大系统的性能，特别是逃逸时序、逃逸分离和逃逸弹道闭环控制的正确性，为我国载人月球探测工程打下了重要的技术基础~\cite{苗志文2026,李喜晟2026}。

\section{本文主要研究内容}

面向载人登月任务需求，本文选取猎户座飞船作为参考构型，开展发射逃逸阶段的六自由度动力学建模、蒙特卡洛仿真与风险评估研究。主要研究内容有：

\begin{enumerate}
    \item 发射阶段逃逸剖面设计。对比载人登月与近地轨道逃逸任务的差异，分析发射阶段零高度和最大动压两种典型逃逸工况的剖面特征。

    \item 关键阶段六自由度动力学仿真模型建立。定义坐标系和变换矩阵，构建国际标准大气（ISA）大气模型、气动力/力矩模型、多发动机联合推进模型和姿控模型，推导完整的平动和转动动力学微分方程组，建立六类八因素干扰项的数学模型。

    \item 打靶仿真与结果分析。设定两种工况的靶点和干扰项分布参数，通过L27正交试验和500组拉丁超立方试验开展批量打靶仿真，统计分析脱靶量、命中率、敏感因素排序和关键指标的分布情况。

    \item 逃逸任务风险性评估。基于前面的打靶统计结果，建立评估指标体系，从命中风险、姿态风险和敏感性三个方面，对两种工况展开评估，找出薄弱环节并提出改进建议。
\end{enumerate}

为更直观地展示本文的研究思路与技术路径，图~\ref{fig:tech-roadmap}~给出了本文的技术路线图。从任务分析、剖面设计出发，经六自由度动力学建模，到蒙特卡洛打靶仿真，最终完成风险评估，形成完整的研究闭环。

\begin{figure}[htbp]
\centering
\includegraphics[width=0.92\textwidth]{技术路线图.png}
\caption{本文技术路线图}
\label{fig:tech-roadmap}
\end{figure}
